% Chapter 3 placeholder
\section{モデル}
本研究では,強相関電子系における電子間の局所的なクーロン相互作用と非局所的な運動エネルギーを記述する,最も基礎的かつ重要な模型であるハバードモデル (Hubbard model)を対象とする.

本研究で扱う単一バンドハバードモデルのハミルトニアンは次のように与えられる.
\begin{equation}
    \hat{H} = -t \sum_{\langle i,j \rangle, \sigma} (\hat{c}_{i\sigma}^\dagger \hat{c}_{j\sigma} + H.c.)+ U \sum_i \hat{n}_{i\uparrow} \hat{n}_{i\downarrow}.
\end{equation}
ここで,第1項の $t$ は最隣接サイト間のホッピング積分であり,第2項の $U$ は同一サイトに2つの電子が占有した (二重占有 )際に生じるクーロン斥力エネルギーである.

本計算においては,格子構造の違いによるモット転移への影響を調べるため,以下の2つの格子モデルを想定してシミュレーションを行う.

\subsection{無限次元Bethe格子}
ハバードモデルの格子が無限次元 (配位数が無限大 )である極限を仮定する.

この無限次元極限のBethe格子において,自由電子の状態密度は以下の半円状態密度で記述される.
\begin{equation}
    D(\epsilon)= \frac{2}{\pi D} \sqrt{1 - \left(\frac{\epsilon}{D}\right)^2} \quad (|\epsilon| \le D).
\end{equation}
ここで $D$ はバンドの半幅であり,本研究ではこのBethe格子におけるエネルギーの基準スケールとして $D=1.0$ を用いる.

\subsection{2次元正方格子}
有限次元の格子系として,2次元正方格子 (2D square lattice)を想定する.運動量空間における自由電子の分散関係は次式で与えられる.
\begin{equation}
    \epsilon_{\mathbf{k}} = -2t (\cos k_x + \cos k_y).
\end{equation}
本研究ではホッピング積分を $t = 0.25$ と設定した.これによりバンド幅は $W = 8t = 2.0$,半バンド幅は $D = 1.0$ となる.

\section{記法とパラメータ}
本論文における ghost Gutzwiller 近似 (gGA)を用いた数値計算は,以下の記法とパラメータ設定に基づいて実行した.

\begin{itemize}
    \item \textbf{軌道数とゴースト自由度}: 
    各サイトの物理的な空間軌道数は $\nu_i = 1$ (スピン自由度を含めて2モード )とした.
    
    変分空間を拡張するgGAの補助空間を構成するため,2つの空間ゴースト軌道 (スピン自由度を含めて4モード )を導入した.
    
    したがって,補助空間全体の空間軌道数は物理軌道とゴースト軌道を合わせて3軌道となり,gGAの拡張パラメータとしては $B=3$ に相当する.
    
    \item \textbf{電子数(充填率)}:
    系の電子数は,各サイト・各スピンあたり $\langle \hat{n}_\sigma \rangle = 0.5$ となる半充填(ハーフフィリング,$n=1$ )状態を対象とした.
    
    \item \textbf{相互作用パラメータ ($U$)とスイープ方向}: 
    金属からモット絶縁体への相転移を詳細に観測するため,クーロン相互作用 $U$ を $0.1$ から $5.0$ の範囲で $0.1$ 刻みで変化させた.
    
    さらに,一次相転移に伴うヒステリシス (金属相と絶縁体相の共存領域 )を捉えるため,以下の2つの経路で計算を行った.
    \begin{enumerate}
        \item \textbf{昇順スイープ (金属 $\to$ 絶縁体 )}: 小さい $U$ 側の金属状態の解を初期値とし,断熱的に $U$ を増加させる.
        \item \textbf{降順スイープ (絶縁体 $\to$ 金属 )}: 大きい $U$ 側のモット絶縁体状態の解を初期値とし,断熱的に $U$ を減少させる.
    \end{enumerate}
    
    \item \textbf{比較のための参照値}:
    gGAの精度を評価するため,無限次元極限において厳密となる動的平均場理論 (DMFT)によって得られているBethe格子の転移点の臨界値 $U_{c1} = 2.39$ (絶縁体解が消失する点 )および $U_{c2} = 2.94$ (金属解が消失する点 )を比較の基準として用いる.
    
    \item \textbf{計算温度 ($T$)}:
    本研究は基底状態の性質を調べる目的であるが,数値アルゴリズムの安定性と最適化の収束性を確保するため,フェルミ分布関数に十分低い人工的な温度 ($T = 0.003$ )を導入して計算を行った.
\end{itemize}